\section{The structural theorem}
We retain the classical rank-three rotation-group correspondence:
a finite generating pair $(a,b)$ of orders at least three satisfying
$(ab)^2=1$ and $\gen a\cap\gen b=1$ determines a directly regular
or chiral polyhedron. It is chiral exactly when no automorphism
simultaneously inverts $a,b$. In the chiral case the full group is
$G=\gen{a,b}$, and tightness means $|G|=|a||b|$.
This statement is the exact geometric interface used in the source's
Section 2.2, attributed there to Schulte and Weiss; it is not newly
proved by the following group calculations.

\begin{theorem}[Source Theorem 1.1]
Let $(G;a,b)$ be the full group of a chiral polyhedron and
$|G|=2p^m$ for an odd prime $p$. After exchanging $a,b$ if necessary,
\[
|a|=p^u,\quad |b|=2p^v,\quad u,v\geq1,
\qquad G=P\rtimes\gen{ab},\quad P=\gen{a,b^2}.
\]
The Sylow subgroup $P$ is nonabelian and needs exactly two generators.
The polyhedron is tight if and only if $P$ is metacyclic.
\end{theorem}
\begin{proof}
The Sylow subgroup $P$ has index two and is normal. The product
$ab$ is a nonidentity involution: equality to one would contradict
the nontrivial disjoint cyclic subgroups. Thus $\gen{ab}$ meets
$P$ trivially, giving the displayed semidirect product.
Both generators cannot be in $P$, because they generate $G$.
Neither can both be outside it, because their product would be
an involution in the odd-order group $P$. Thus exactly one is outside;
orders and nondegeneracy give the displayed type.
Put $x=a$ and $y=b^2$. The identities
\[
x^b=y^{-1}x^{-1},\qquad y^b=y
\]
show $\gen{x,y}$ is normalized by $a,b$. Its quotient in $G$
is generated by the image of $b$, of order at most two.
It is contained in $P$, so equals $P$. If $P$ were cyclic, the
two nontrivial cyclic subgroups $\gen x,\gen y$ would share
its unique subgroup of order $p$, contradicting the intersection
condition. Hence two generators are necessary as well as sufficient.
If $P$ were abelian, the abelian-kernel reversal lemma of~\cite{UF} with $t=ab$ would
supply a reversal, contradicting chirality. In particular $m\geq3$.

For the last assertion we retain precisely two classical odd-primary
subgroup facts, also invoked as source Theorems 2.4 and 2.5:
(i) a finite odd $p$-group which is a setwise product of two cyclic
subgroups is metacyclic (Huppert, \emph{Endliche Gruppen I},
Theorem 11.5 as cited in the source);
(ii) in a finite metacyclic odd $p$-group, subgroups permute
(the Iwasawa result cited through Schmidt, \emph{Subgroup Lattices
of Groups}, Theorem 2.3.1).
These are retained classical inputs with full proof cost unknown,
not deductions from the finite diagnostics.
The two cyclic subgroups here intersect trivially. If the polyhedron
is tight, their setwise product has size $p^{u+v}=p^m=|P|$;
it is $P$, so (i) applies. Conversely (ii) makes their setwise
product a subgroup containing the two generators of $P$.
It is therefore $P$, and the order formula gives tightness.
\end{proof}

